\section{Three Recurring Failure Patterns}
\label{sec:failure}

Across the four axes surveyed in Sections~\ref{sec:geometry}--\ref{sec:sufficiency},
and through our own experimental programme, we identify three structural reasons
why importing a modern mathematical framework into granular computing's core
pipeline tends not to produce a publishable contribution---or produces one that
does not stand up to independent scrutiny.  We call these \emph{Pattern~A}
(reduction to a known quantity), \emph{Pattern~B} (dominance of
label-orthogonal geometry),
and \emph{Pattern~C} (self-optimising evaluation); Figure~\ref{fig:patterns}
maps each pattern to the stage of the GrC pipeline where it arises.
They are not mutually exclusive; a single proposal can exhibit all three.
Section~\ref{sec:trends} shows how three current research directions
(incremental selection, three-way decision, neural-augmented GrC) can
navigate these patterns when evaluated correctly.

%-----------------------------------------------------------------------
\subsection{Pattern A: Reduction to a Known Granular Quantity}
\label{sec:patA}
%-----------------------------------------------------------------------

The most common failure is definitional: once the new quantity is unfolded,
it coincides with---or is bounded trivially by---a quantity already present
in the granular computing literature.

\paragraph{Example~A1: H$_0$ persistence and single-linkage clustering.}
The zero-dimensional persistent homology diagram of a point cloud under a
Vietoris--Rips filtration records the birth and death of connected components.
The death of a component at scale $\varepsilon$ marks the moment two
components merge---exactly the linkage event in single-linkage hierarchical
clustering.  This correspondence is a standard result: the $H_0$ barcode
of a Vietoris--Rips filtration is the merge tree of single-linkage
clustering, equivalently the minimum spanning tree's edge
weights~\cite{CarlssonMemoli2010hierarchical,Chazal2013persistenceclustering}.
Consequently, any attribute-importance score derived from the
H$_0$ diagram on the full feature set encodes the same information as the
single-linkage dendrogram, a quantity studied since the 1960s and available
in $O(n^2)$ time.  Example~A1 is therefore a genuine reduction, not an
analogy: the two quantities are equal, not merely correlated.

\paragraph{Example~A2: Supervised separation persistence and positive region.}
Consider a nearest-enemy superlevel filtration: each object $x$ is
``safe'' at radius $t$ if its nearest enemy is at distance $\geq t$;
connected components of safe objects form $H_0$ features as $t$ decreases.
Let $\phi(B)$ be the total persistence of these components under $B$.
An attribute is important if removing it reduces $\phi$.  After algebraic
simplification, the dominant term in $\phi(B)$ tracks the radius at which
class cores first connect to an enemy---that is, $1 - \DEP(B)$, where
$\DEP(B) = |\POS_B(d)|/|\U|$ is the NRS dependency degree at the same
radius.  The residual terms encode within-class boundary complexity, but
in the standard NRS regime where class cores are the dominant signal,
supervised persistence and NRS dependency convey the same ranking
information.  We state this as a heuristic reduction rather than an
identity: unlike Example~A1 we have no proof that the residual terms are
uninformative, only the observation that they vanish when every class core
is separated from its nearest enemy by a single scale.  The weaker claim is
enough for the diagnostic purpose---a proposal in this shape carries the
burden of showing that its residual is not negligible---but it is not
enough to assert equality, and we do not.

\paragraph{Example~A3: Fuzzy-T2 discernibility profiles and maximal
discernibility pairs.}
The proposal to use a fuzzy T$_2$ (Hausdorff-separation) degree between
attribute \emph{profiles}---vectors $\Delta_a(i,i') = |a(x_i)-a(x_{i'})|$
over cross-class pairs---as a basis for partitioning attributes shares its
core insight with Dai et al.~\cite{Dai2018TFS}: algorithms that operate on
object pairs can ignore already-discerned pairs, reducing cost from $O(m \cdot n^2)$
to $O(m\cdot|D|)$ where $D$ is the active discernibility set.
The T$_2$ framing adds topological language without adding
computational or theoretical content beyond that 176-cited result.

\paragraph{Diagnostic.}
Pattern~A is diagnosable before any experiment: write out the new quantity
symbolically and ask whether it is a monotone transformation of an existing
granular measure.  If the Jacobian of the mapping is bounded and non-vanishing,
the two quantities carry the same information up to a scaling factor.

%-----------------------------------------------------------------------
\subsection{Pattern B: Dominance of Label-Orthogonal Geometry}
\label{sec:patB}
%-----------------------------------------------------------------------

The second failure arises when the new importance signal is dominated by
properties of a feature's distribution that carry no information about the
decision task.  Measurement scale is the most familiar such property and the
one a stability bound like Theorem~\ref{thm:pai_bound} makes explicit, so we
begin there---but, as the evidence below shows, it is a special case rather
than the whole pattern, and a proposal can fail this way while surviving a
scale check.

\paragraph{Experimental evidence.}
We computed persistence-based attribute importance (PAI)---bottleneck distance
between the full persistence diagram and the diagram after removing attribute
$j$---on ten UCI datasets.  Each of five seeds draws one 80\% subsample without
replacement, giving five importance vectors per method and dataset
(total: $4{,}975$ measurements; raw data released as
\texttt{raw\_importance.csv}).  Downstream accuracies are separate: they come
from stratified 5-fold cross-validation repeated over the same five seeds,
with every importance ranking recomputed inside each training fold
(Section~\ref{sec:protocol}).
Table~\ref{tab:iris_top} reports the top-ranked attribute by PAI and by NRS
dependency on the Iris dataset as a representative example.

% AUTO-GENERATED by gen_tables.py -- do not edit by hand.
% Regenerate:  python3 gen_tables.py   Audit:  python3 verify_tables.py
\begin{table}[pos=!htbp]
\caption{Top-ranked attribute by method on Iris ($n=150$, $d=4$).  The
top-1 attribute is the arg-max of mean importance over the five seed
subsamples; the accuracies are mean $\pm$ SD over the 25 cross-validation
runs (5 folds $\times$ 5 seeds) using only that attribute. PAI
selects the attribute whose removal most disrupts the persistence diagram;
NRS selects the most discriminative attribute for the class label.}
\label{tab:iris_top}
\centering\small
\begin{tabular}{lccc}
\toprule
Method & Top-1 attribute & 3-NN acc.\ (top-1) & SVM acc.\ (top-1) \\
\midrule
PAI  & \texttt{sepal\_width}  & $48.9 \pm 7.4$\% & $54.0 \pm 8.4$\% \\
NRS  & \texttt{petal\_width}  & $95.9 \pm 3.9$\% & $96.0 \pm 3.5$\% \\
\bottomrule
\end{tabular}
\end{table}


PAI ranks \texttt{sepal\_width} first because removing it causes the largest
bottleneck-distance shift in the persistence diagram, even though it
contributes no class signal at the neighbourhood scale used by NRS.
\texttt{petal\_width}, the most discriminative attribute (NRS top-1),
receives a smaller PAI score because the persistence diagram is less
sensitive to its removal.
The downstream consequence is stark: selecting on PAI yields 3-NN accuracy of
48.9\% (barely above the 33\% three-class chance level) versus 95.9\%
for NRS.  The same pattern holds under SVM-RBF evaluation: PAI top-1 gives
54.0\% versus NRS top-1 at 96.0\%.

\paragraph{Generalisation across datasets.}
The Iris result is not an artefact of one dataset.
Table~\ref{tab:rho} reports the Spearman rank correlation $\rho$ between PAI
and NRS importance across all ten benchmarks.

% AUTO-GENERATED by gen_tables.py -- do not edit by hand.
% Regenerate:  python3 gen_tables.py   Audit:  python3 verify_tables.py
\begin{table}[pos=!htbp]
\caption{Spearman $\rho$ between PAI and NRS importance rankings per dataset
(mean importance over 5 seeds before ranking; $d$ = number of attributes,
which is also the sample size for the Spearman test).
Note: datasets with $d \leq 7$ cannot reach $p < 0.05$ at any $\rho$.}
\label{tab:rho}
\centering\small
\begin{tabular}{lccc}
\toprule
Dataset & $d$ & $\rho$ & $p$-value \\
\midrule
Iris       &  4 & $-0.80$ & $0.20$ \\
Wine       & 13 & $-0.29$ & $0.34$ \\
Ionosphere & 34 & $-0.17$ & $0.35$ \\
Sonar      & 60 & $+0.04$ & $0.76$ \\
Glass      &  9 & $+0.06$ & $0.88$ \\
Ecoli      &  7 & $+0.71$ & $0.07$ \\
Heart      & 13 & $+0.08$ & $0.79$ \\
Parkinsons & 22 & $+0.30$ & $0.17$ \\
Seeds      &  7 & $+0.21$ & $0.64$ \\
WDBC       & 30 & $-0.50$ & $0.01$ \\
\midrule
Mean (10)  & --- & $-0.034$ & --- \\
\bottomrule
\end{tabular}
\end{table}


Table~\ref{tab:rho} shows the ten-dataset picture.
The three-dataset pilot (Iris, Wine, Ionosphere) produced $\bar\rho=-0.42$,
suggesting complementarity.  The full run gives $\bar\rho = -0.034$, consistent
with zero correlation.  A Fisher-$z$ weighted aggregate (weights $d_i - 3$)
gives $\bar\rho_z = -0.060$, 95\% CI $[-0.207,\,+0.091]$, with
heterogeneity $Q = 15.8$, $I^2 = 43\%$---the confidence interval straddles
zero, confirming no systematic relationship.
Only one $\rho$ is individually significant
(WDBC, $d=30$, $\rho=-0.50$, $p=0.01$); on datasets with $d \leq 7$
(Iris, Ecoli, Seeds), the Spearman test has no power to reject the null
at any ranking.  Six rhos are positive, consistent with PAI and NRS ranking
attributes in the same order on more datasets than not, but a sign test on
6/10 is not significant ($p = 0.75$, two-sided binomial): the two methods are
uncorrelated, not complementary.  The pilot's apparent complementarity was not representative.

\paragraph{Classifier robustness (SVM-RBF).}
Following our own recommended practice (Section~\ref{sec:standard}), we
repeat the analysis with SVM-RBF ($C=1$, $\gamma=\text{scale}$), a
classifier whose geometry differs from the $k$-NN rule.  Over $k < d$,
mean SVM accuracy is: NRS $85.1\%$, PAI $80.2\%$ ($p = 0.006$, dataset-level
Wilcoxon, $n=10$; PAI wins on 1/10 datasets).  The SVM results confirm
the 3-NN finding: PAI's inferiority is not an artefact of the classifier.

PAI is also expensive.  It is $\sim 380\times$ slower than NRS: mean of per-dataset median reduct time
is 6.5\,s for PAI vs 0.017\,s for NRS (on the same single-threaded harness).
ReliefF (0.068\,s) and mutual information (0.014\,s) are also orders of
magnitude faster.  One anomaly: Sonar ($d=60$) reports PAI median time of
0.018\,s, comparable to NRS (0.034\,s), despite having the highest
dimensionality; this suggests a computation shortcut or caching effect at
that particular filtration setting that warrants further investigation.
The general cost is inherent: PAI requires a full Rips
persistence computation per attribute per subsample, with $O(n^3)$ worst-case
complexity.

\paragraph{Additional baselines: hybrid and null.}
The released data includes a hybrid ranking (\texttt{pai\_nrs}: the
element-wise mean of z-scored PAI and NRS importance vectors) and a variance
ranking (features ordered by sample variance on the z-scored data---effectively
a null baseline, since all variances centre near 1.0).  Neither rescues PAI.
Over $k < d$, mean 3-NN accuracy is: NRS $83.7\%$, PAI+NRS hybrid $80.1\%$
($p = 0.010$ vs NRS, Wilcoxon), PAI $77.6\%$ ($p = 0.004$ vs PAI+NRS).
The variance baseline achieves $79.0\%$---statistically indistinguishable
from PAI ($p = 0.23$).  That a null ranking matches PAI is the sharpest
available statement of the negative result.

\paragraph{Root cause: theoretical bound vs empirical behaviour.}
Theorem~\ref{thm:pai_bound} guarantees $\mathrm{PAI}(j) \leq \mathrm{range}(a_j)$:
the bound is proportional to the feature's range.  One might therefore
hypothesise that PAI scores track range, making them scale artefacts.
We test this directly.  Table~\ref{tab:pai_range} reports the Spearman
$\rho$ between PAI and $\mathrm{range}(a_j)$ per dataset on the z-scored
data actually used.

% AUTO-GENERATED by gen_tables.py -- do not edit by hand.
% Regenerate:  python3 gen_tables.py   Audit:  python3 verify_tables.py
\begin{table}[pos=!htbp]
\caption{Spearman $\rho$ between mean PAI score and
$\mathrm{range}(a_j)$ per dataset.  Data is z-score standardised,
so range variation across features reflects kurtosis and outlier
structure, not raw measurement scale.}
\label{tab:pai_range}
\centering\small
\begin{tabular}{lccc}
\toprule
Dataset & $d$ & $\rho(\text{PAI},\text{range})$ & $p$ \\
\midrule
Iris       &  4 & $+0.80$ & $0.200$ \\
Wine       & 13 & $+0.64$ & $0.019$ \\
Ionosphere & 34 & $+0.42$ & $0.014$ \\
Sonar      & 60 & $+0.33$ & $0.010$ \\
Glass      &  9 & $-0.52$ & $0.154$ \\
Ecoli      &  7 & $-0.54$ & $0.215$ \\
Heart      & 13 & $-0.04$ & $0.887$ \\
Parkinsons & 22 & $-0.85$ & $<0.001$ \\
Seeds      &  7 & $+0.89$ & $0.007$ \\
WDBC       & 30 & $+0.35$ & $0.060$ \\
\midrule
Mean (10)  & --- & $+0.15$ & --- \\
\bottomrule
\end{tabular}
\end{table}


The mean correlation is $+0.15$ and highly heterogeneous: Parkinsons is
\emph{significantly negative} ($\rho = -0.85$, $p < 0.001$).  The
theoretical bound of Theorem~\ref{thm:pai_bound} is far from tight
(mean ratio $0.038$, max $0.194$; equivalently $0.075$ and $0.387$ against
the sharper per-dimension bound), so range is not the \emph{primary} driver
of PAI on most datasets.  The
mechanism is subtler: PAI measures how much removing a coordinate
changes the \emph{global topology} of the point cloud, a quantity that
depends on the interaction between the coordinate's distribution (kurtosis,
outlier structure after z-scoring) and the dataset's geometric embedding.
This sensitivity to global geometry rather than local class structure is why
PAI fails as a supervised importance measure.

This is also why we define Pattern~B by label-orthogonality rather than by
scale.  Had we kept the narrower definition, our own flagship example would
not satisfy it: PAI's correlation with range is significantly positive on
only four of ten datasets and significantly \emph{negative} on one, so
``dominated by range'' is not what the data show.  What the data do show is
that PAI is computed entirely from the geometry of the unlabelled point
cloud, is uncorrelated with discriminability (Table~\ref{tab:rho}), and is
indistinguishable from a variance null baseline downstream ($p=0.23$).  The
failure is that the imported quantity measures something real about the data
that has no bearing on the decision; measurement scale is the instance of
that failure a stability theorem can bound, which is why it is the one the
literature notices.

Note that the data is z-score standardised (Section~\ref{sec:protocol}), so
all features have unit variance; the residual range variation comes from
kurtosis and outliers.  This makes the negative result stronger
than it would be on raw data: even after removing the obvious scale artefact
via standardisation, PAI still does not align with discriminative importance.

\paragraph{Generality.}
Pattern~B is not specific to persistent homology.  Any importance measure
$s_j$ that is computed from the geometry of the feature distribution without
reference to the label is at risk, and a bound depending on measurement-scale
properties of $a_j$ (range, variance, inter-quantile spread) is the visible
symptom rather than the disease.  For purely
distance-based importance (e.g.\ Relief, HSIC with Gaussian kernel), the
community has learned to z-score inputs before applying such methods.  The
PAI case shows that z-scoring is necessary but not sufficient: homological
stability introduces a dependency on the \emph{shape} of each coordinate's
distribution that standardisation does not remove.

%-----------------------------------------------------------------------
\subsection{Pattern C: Self-Optimising Evaluation}
\label{sec:patC}
%-----------------------------------------------------------------------

The third failure is evaluative: the new method is benchmarked against metrics
that are algebraically related to---or optimised by construction by---the
method's own objective.

\paragraph{Experimental evidence.}
We evaluated five feature-selection methods (PAI, NRS,
ReliefF~\cite{Kononenko1994,RobnikKononenko2003}, mutual information
(mRMR~\cite{Peng2005mRMR}), and variance) on structure-preservation metrics: Wasserstein distance
of persistence diagrams (\texttt{wass\_persist}), trustworthiness
(\texttt{trust}), $k$-NN overlap (\texttt{knn\_overlap}), and adjusted Rand
index of cluster structure (\texttt{cluster\_ari}).
Results across ten datasets are in Table~\ref{tab:structure}.

% AUTO-GENERATED by gen_tables.py -- do not edit by hand.
% Regenerate:  python3 gen_tables.py   Audit:  python3 verify_tables.py
\begin{table}[pos=!htbp]
\caption{Structure-preservation comparison, PAI vs NRS (10 datasets).
\texttt{wass\_persist} is lower-is-better (PAI's own metric); the others are
higher-is-better and are not algebraically tied to PAI's objective.  Note that
\texttt{trust} and \texttt{knn\_overlap} share a neighbourhood-size parameter,
fixed here at the conventional $K=10$ and ablated in
Table~\ref{tab:knn_scale}.  Means, SDs, win counts
and tests all use the same aggregation: one value per dataset, averaged over
$k<d$.  $p$-values: dataset-level paired Wilcoxon signed-rank
($n=10$, two-sided).}
\label{tab:structure}
\centering\small
\begin{tabular}{lcccc}
\toprule
Metric & PAI (mean $\pm$ SD) & NRS (mean $\pm$ SD) & PAI wins / 10 & $p$ \\
\midrule
\texttt{wass\_persist}  & $66.7 \pm 40.1$ & $77.4 \pm 51.2$ & 8/10 & $0.049$ \\
\texttt{knn\_overlap}   & $0.530 \pm 0.075$ & $0.474 \pm 0.090$ & 8/10 & $0.064$ \\
\texttt{trust}          & $0.889 \pm 0.024$ & $0.903 \pm 0.030$ & 1/10 & $0.006$ \\
\texttt{cluster\_ari}   & $0.482 \pm 0.122$ & $0.665 \pm 0.119$ & 1/10 & $0.004$ \\
\bottomrule
\end{tabular}
\end{table}


PAI wins on \texttt{wass\_persist} (8/10 datasets, $p=0.049$; lower is
better)---because PAI selects features that best preserve the persistence
diagram, precisely what Wasserstein distance measures.  This is the only
metric of the four on which PAI's advantage reaches significance, and it is
the one algebraically tied to PAI's objective.  PAI also scores higher on
\texttt{knn\_overlap} on 8/10 datasets, at $p=0.064$ under the default
neighbourhood size $K=10$.  Rather than explain this away, we tested it.

\paragraph{Is the \texttt{knn\_overlap} edge a scale artefact?  A direct test.}
The natural explanation is that PAI's bottleneck objective and $k$-NN overlap
probe the same local neighbourhoods, so the advantage should be largest when
$K$ matches the scale PAI's filtration actually sees and should decay as the
two are mismatched.  That explanation is post-hoc until tested, so we ran the
discriminating experiment: holding the selected feature subsets fixed at those
of Table~\ref{tab:structure}, we recomputed the metric for
$K \in \{3,5,10,20,30,50\}$.  Crucially we did the same for
\texttt{trust}, which takes the same neighbourhood parameter and is
also conventionally left at $K=10$---ablating only the metric that favours
PAI would repeat the selective reporting this survey criticises.
Table~\ref{tab:knn_scale} reports both.

% AUTO-GENERATED by gen_tables.py -- do not edit by hand.
% Regenerate:  python3 gen_tables.py   Audit:  python3 verify_tables.py
\begin{table}[pos=!htbp]
\caption{Ablation of the neighbourhood parameter $K$ shared by the two
local structure metrics.  Feature subsets are exactly those of
Table~\ref{tab:structure}; only $K$ changes.  Each row aggregates over $k<d$
per dataset, then compares PAI with NRS by dataset-level paired Wilcoxon
signed-rank ($n=10$, two-sided); $\Delta = \mathrm{PAI}-\mathrm{NRS}$, and
$^{*}$ marks $p<0.05$.  $K=10$ is the default used in
Table~\ref{tab:structure} and throughout the literature.}
\label{tab:knn_scale}
\centering\small
\begin{tabular}{cccccc}
\toprule
$K$ & PAI & NRS & $\Delta$ & PAI wins / 10 & $p$ \\
\midrule
\multicolumn{6}{l}{\emph{\texttt{knn\_overlap}} (higher is better)} \\
 3 & $0.407$ & $0.328$ & $+0.079$ & 8/10 & $0.037$\,$^{*}$ \\
 5 & $0.459$ & $0.387$ & $+0.072$ & 8/10 & $0.037$\,$^{*}$ \\
10 & $0.530$ & $0.474$ & $+0.056$ & 8/10 & $0.064$ \\
20 & $0.599$ & $0.570$ & $+0.029$ & 7/10 & $0.160$ \\
30 & $0.643$ & $0.628$ & $+0.014$ & 7/10 & $0.275$ \\
50 & $0.697$ & $0.710$ & $-0.013$ & 2/10 & $0.193$ \\
\addlinespace
\multicolumn{6}{l}{\emph{\texttt{trust}} (higher is better)} \\
 3 & $0.891$ & $0.904$ & $-0.012$ & 1/10 & $0.004$\,$^{*}$ \\
 5 & $0.890$ & $0.903$ & $-0.013$ & 1/10 & $0.004$\,$^{*}$ \\
10 & $0.889$ & $0.903$ & $-0.014$ & 1/10 & $0.006$\,$^{*}$ \\
20 & $0.889$ & $0.906$ & $-0.017$ & 1/10 & $0.004$\,$^{*}$ \\
30 & $0.890$ & $0.910$ & $-0.020$ & 1/10 & $0.004$\,$^{*}$ \\
50 & $0.890$ & $0.914$ & $-0.023$ & 1/10 & $0.004$\,$^{*}$ \\
\bottomrule
\end{tabular}
\end{table}


The advantage is entirely a small-$K$ phenomenon.  It is significant at
$K=3$ and $K=5$ ($p=0.037$), borderline at the default $K=10$ ($p=0.064$),
absent by $K=20$ ($p=0.16$), and at $K=50$ it reverses: PAI wins on
2/10 datasets and trails NRS by $0.013$.  An analyst who reports
\texttt{knn\_overlap} at $K=5$ concludes that PAI preserves neighbourhood
structure significantly better than NRS; an analyst who reports the same
metric at $K=50$ concludes the opposite, from the same feature subsets.  The
conclusion is determined by a free parameter of the evaluation metric.

The mechanism, however, is not the one we predicted.  If the advantage came
from a match between $K$ and PAI's own filtration scale, the $K$ maximising
the advantage would track that scale across datasets.  It does not.  Taking
the empirical filtration scale of each dataset to be the mean number of
points within the Rips threshold ($\varepsilon=2$ on standardised data), which
ranges from $1.6$ on Heart to $77.1$ on Ecoli, the rank correlation between
that scale and the advantage-maximising $K$ is $\rho=-0.10$ ($p=0.80$,
$n=9$ datasets with a non-degenerate scale); pooling all dataset--$K$ pairs,
the correlation between scale mismatch $|\log K - \log k_{\mathrm{PAI}}|$ and
the advantage is $\rho=-0.05$ ($p=0.74$).  Seven of ten datasets peak at
$K=3$ whether their filtration scale is $2$ or $77$.  We therefore withdraw
the same-neighbourhoods explanation: the edge is a generic small-$K$ effect,
not a signature of PAI's geometry.
(One dataset is degenerate in a way worth recording: on Sonar, $d=60$, no
point has \emph{any} neighbour within $\varepsilon=2$ after standardisation,
so PAI's filtration sees no connectivity at all---and Sonar is the one
dataset where PAI trails NRS at every $K$.)

The comparison with \texttt{trust} is what makes this conclusive.  That
metric carries the identical parameter, so if PAI's verdicts were generally
$K$-sensitive we should see it move too.  It does not move at all: PAI loses
on 1/10 datasets at every $K$ from 3 to 50, at $p \leq 0.006$ throughout, and
the gap widens with $K$ (from $-0.012$ at $K=3$ to $-0.023$ at
$K=50$).  The third metric, \texttt{cluster\_ari}, has no free local scale at
all---its cluster count is fixed to the number of classes---and agrees
($p=0.004$).

The asymmetry is the finding.  The metric on which PAI appears to lead is
precisely the one whose verdict an analyst selects when choosing $K$; the
metrics on which PAI loses return the same verdict at every setting of that
same parameter.  The negative result is therefore robust to the choice we
have just shown to be decisive elsewhere, and the apparent positive result is
not.  A paper that evaluates a
persistence-based selector only on persistence-based (or same-scale) downstream
metrics has demonstrated self-consistency, not generalisation.

\paragraph{Generality.}
Pattern~C recurs whenever evaluation and design share the same mathematical
object.  A coverage-based reduct evaluated by coverage, a mutual-information
selector evaluated by mutual information with the label, a graph-regularised
model evaluated by graph smoothness---all exhibit the same circularity.
The diagnostic is simple: if removing the proposed evaluation metric from the
paper would remove the only result where the method wins, Pattern~C is present.
A second form is subtler and, on the evidence of Table~\ref{tab:knn_scale}, at
least as common: if the metric carries a free parameter and the method's win
survives only in part of that parameter's range, the win is a property of the
parameter choice, not of the method.  Report the whole range, and check
whether the metrics that \emph{disagree} with the method are equally
sensitive---if they are not, that asymmetry is the result.

\paragraph{Recommended practice.}
For any proposed importance measure $s$, include at least one evaluation
metric $m$ such that $s \not\to \arg\max m$: that is, optimising $s$ does not
provably optimise $m$.  Downstream accuracy on held-out data with a classifier
that does not use $s$'s geometry (e.g., SVM-RBF or Random
Forest~\cite{Breiman2001RF} when $s$ is geometry-based)
serves this purpose if statistical significance is reported.

%-----------------------------------------------------------------------
\subsection{Diagnostic Checklist}
\label{sec:checklist}
%-----------------------------------------------------------------------

Before investing in implementation, a proposed GrC extension should pass
three checks:

\begin{enumerate}
\item \textbf{Anti-A}: Write the new quantity symbolically.  Is it a
monotone function of an existing granular quantity?  If yes, prove it is
strictly more informative, or abandon.
\item \textbf{Anti-B}: Is the quantity computed from the feature
distribution without reference to the label?  If so, standardisation is the
first test but not a sufficient one---our own experiments were run on
z-scored data and Pattern~B appeared regardless.  Run two checks, not one.
(i)~Correlate the score against a pure scale statistic (range, IQR) on the
standardised data; a significant positive correlation is a failure.
(ii)~More important, correlate it against an established
\emph{label-aware} importance measure: a score that is uncorrelated with
discriminability has failed even if it passes~(i), which is precisely what
PAI does (Table~\ref{tab:rho} versus Table~\ref{tab:pai_range}).  A
label-free geometric score must earn its place by predicting the label
better than a null ranking, not merely by being unlike range.
\item \textbf{Anti-C}: Identify at least one evaluation metric that is
not algebraically related to the method's objective.  Report that
metric as the primary result.
\end{enumerate}

These checks cost one afternoon; their absence costs months of wasted
implementation.
